\section{Methodology}
\label{sec:method}

The pipeline consists of five steps. I detail each below, pointing out the
formulas that define its geometry.

\subsection{Perspective Rectification}
\label{sec:method:rectify}
I first warp the phone image to a fronto-parallel view. I binarize the page from
the background using Otsu's method, which finds the threshold $t^\star$ that
maximizes the between-class variance
\begin{equation}
  t^\star = \arg\max_{t}\; \omega_0(t)\,\omega_1(t)\,\big[\mu_0(t)-\mu_1(t)\big]^2,
  \label{eq:otsu}
\end{equation}
where $\omega_i$ and $\mu_i$ represent the probability mass and mean of the two
pixel classes separated by $t$~\cite{otsu1979threshold}. The four corners of the
page are extracted from the page mask via a cascade of increasingly liberal
detectors: intersections between Hough lines~\cite{duda1972hough}, polygonal
approximation of the contours~\cite{suzuki1985contours}, and the minimum-area
bounding rectangle. The first of these to result in a sufficiently large convex
quadrilateral is selected. They determine a $3\times3$ homography $\mathbf{H}$
that transforms any source point $\mathbf{x}=(x,y,1)^\top$ into the rectified
space as
\begin{equation}
  \tilde{\mathbf{x}}' \simeq \mathbf{H}\,\tilde{\mathbf{x}},
  \qquad
  x' = \frac{h_{11}x+h_{12}y+h_{13}}{h_{31}x+h_{32}y+h_{33}},
  \label{eq:homography}
\end{equation}
where $y'$ is defined analogously~\cite{hartley2004multiple}. I check that a
rectification is valid by computing the residual slant of staff lines in the
warped image and by counting the fraction of background pixels which overlap with
the border of the page.

\subsection{Staff Detection}
\label{sec:method:staves}
On the rectified page I use adaptive thresholding and then the Hough line
transform, which allows each edge pixel to vote for the lines that pass through it
\begin{equation}
  \rho = x\cos\theta + y\sin\theta .
  \label{eq:hough}
\end{equation}
Horizontal line candidates ($\theta\approx 90^{\circ}$) are grouped by their
$y$-intercept and combined to form five-line staves. A group of five line
candidates is a stave if the spacing $d_k$ between the $k$-th and $(k{+}1)$-th
line does not vary too much, measured by the coefficient of variation
\begin{equation}
  \mathrm{CV} = \frac{\sigma_d}{\mu_d} < 0.4 .
  \label{eq:cv}
\end{equation}
The average spacing, $s=\mu_d$, provides the fundamental length scale for all
subsequent steps.

\subsection{Staff Removal and Segmentation}
\label{sec:method:segment}
The staff lines are removed by suppressing thin horizontal runs. The gaps this
creates in the note stems are bridged using a vertical morphological closing of
height $\approx s$. Connected components are then extracted as
symbols~\cite{suzuki1985contours}, and rejected if their area, relative to $s^2$,
indicates they are noise.

\subsection{Symbol Classification}
\label{sec:method:cnn}
The $64\times64$ binary images of symbols are passed through \textsc{SymbolCNN},
a four-layer convolution--batchnorm--pooling CNN followed by two fully connected
layers. This is trained using $35{,}933$ crops from PrIMuS~\cite{calvo2018primus},
labelled with nine possible symbols: whole, half, quarter, and eighth notes; four
different rest types; and \emph{other}. Since the classes are not balanced, I
minimize a weighted cross-entropy loss
\begin{equation}
  \mathcal{L} = -\sum_{c} w_c\, y_c \log \hat{y}_c,
  \qquad
  w_c \propto \frac{1}{\sqrt{n_c}},
  \label{eq:wce}
\end{equation}
where $n_c$ is the number of crops in class $c$, which weights the less common
classes more heavily. I assign each crop its correct label based on a strict rule
and discard potentially ambiguous symbol--image pairs.

\subsection{Pitch Inference and a Geometry-Only Reader}
\label{sec:method:pitch}
Given the centroid of a notehead located at pixel $y_{\text{head}}$ and the stave
containing it, whose bottom line is at pixel $y_{\text{bottom}}$, I calculate its
position on the stave (in units of half-spaces)
\begin{equation}
  p = \mathrm{round}\!\left(\frac{y_{\text{bottom}}-y_{\text{head}}}{s/2}\right).
  \label{eq:staffpos}
\end{equation}
I map a position on a treble-clef staff (bottom line E4) to a MIDI number by a
deterministic mapping, and provide the associated note duration as determined by
the classifier.

\paragraph{Geometry-only reader.}
Since I find that classification is the least robust step
(Sec.~\ref{sec:results}), I present a reader that performs segmentation-free pitch
inference. After a preliminary filling-in of noteheads using connected-component
analysis on the background pixels, I locate candidates using normalized
cross-correlation with a synthetic ellipse $T$ of size proportional to $s$
\begin{equation}
  R(u,v) = \frac{\sum_{i,j} \big(I_{u,v}(i,j)-\bar I_{u,v}\big)\big(T(i,j)-\bar T\big)}
                {\sqrt{\sum_{i,j}\big(I_{u,v}(i,j)-\bar I_{u,v}\big)^2\;\sum_{i,j}\big(T(i,j)-\bar T\big)^2}},
  \label{eq:ncc}
\end{equation}
where $I_{u,v}$ refers to the image patch at $(u,v)$ beneath the template.
Candidates at peaks of $R$ are found and their neighbors removed with greedy
non-maximum suppression. I discard detections of time signatures or clef margins,
and I filter rests and stems based on their solidity: candidates that remain
solid under a notehead-sized opening are likely true noteheads. Final candidates
are mapped to their MIDI pitch using Eq.~\eqref{eq:staffpos} and staff-position
information.
